\documentclass[10pt,a4paper]{amsart}
\usepackage[margin=19mm]{geometry}
\usepackage{amsmath,amssymb,mathtools}
\usepackage[scr=rsfs,cal=euler]{mathalfa}
\usepackage{xfrac}
\usepackage[unicode,hidelinks,bookmarks=false]{hyperref}
\newcommand{\ie}{i.e.\ }
\newcommand{\cf}{cf.\ }
\newcommand{\resp}{respectively\ }
\newcommand{\Subsection}{\subsection}
\providecommand{\nobreakdash}{\nobreak\hbox{-}}
\setlength{\emergencystretch}{2em}
\multlinegap0pt
\usepackage{bm}
\usepackage{bbm}
\usepackage{dsfont}
\usepackage{xspace}
\usepackage{cancel}
\usepackage{mathtools}
\usepackage{amsthm}
\makeatletter
\@ifundefined{theorem}{\newtheorem{theorem}{Theorem}}{}
\@ifundefined{lemma}{\newtheorem{lemma}[theorem]{Lemma}}{}
\@ifundefined{proposition}{\newtheorem{proposition}[theorem]{Proposition}}{}
\makeatother
\DeclareMathOperator{\cont}{cont}
\newcommand{\NN}{\mathbb{N}}
\newcommand{\cc}{\mathbf{c}}
\renewcommand{\ge}{\geqslant}
\renewcommand{\le}{\leqslant}
\renewcommand{\pmod}[1]{\,(\mathrm{mod}\,#1)}
\title[Average root numbers in two isotrivial families of elliptic ]{Average root numbers in two isotrivial families of elliptic curves}
\author{Yijie Diao}
\date{Source: arXiv:2608.10702v1 (CC BY 4.0)}
\begin{document}
\maketitle

\noindent\emph{Source and licence.} Average root numbers in two isotrivial families of elliptic curves. Version arXiv:2608.10702v1 (CC BY 4.0), distributed under the Creative Commons Attribution 4.0 International licence, as recorded in the arXiv licence metadata for this exact version and shown on the abstract page https://arxiv.org/abs/2608.10702v1. Full rights notice: licenses/arxiv-2608.10702-CC-BY-4.0.txt.\par\smallskip

\noindent\textit{Editorial scope.} Counted unit: the Proposition whose source label is \texttt{prop:squarefull-polynomial-values} in the article (source0.tex lines 1862--1877, source label \texttt{prop:squarefull-polynomial-values}), delivered together with the article's own complete original proof of it (source0.tex lines 1879--1974, one \texttt{proof} environment of 96 lines). No mathematics is written, repaired or re-derived: statement and proof are the article's own text line for line. Required context retained uncounted: lem:squarefull-kernel (lines 566--586); lem:nu-bounds (lines 1710--1718). Every definition, display and label of those lines is retained unchanged. Omitted from this package and not redistributed: 1--565, 587--1709, 1719--1861, 1878--1878, 1975--2801. Nothing inside a retained excerpt is omitted or altered: every retained body is delivered exactly as the article's lines are written, so no omission receipt of any kind accompanies this package. Citations: the retained excerpts carry no citation command at all, so no bibliography entry of the article is transcribed and no reference is redistributed. Reference adaptation: none was needed and none was made; every reference inside the delivered unit resolves inside it, so no unresolved reference or citation is produced. Printed slips kept literal: no printed word, symbol, hypothesis or number of the counted statement or of its proof was altered. Layout and class: the article's own class is \texttt{amsart} with packages including fontenc, amsmath, amssymb, amsfonts, amsthm, geometry, etoolbox, hyperref; the wrapper is the lane's pinned \texttt{amsart} editorial wrapper at 10pt on A4, which restates the theorem-like environments the retained text needs and adds the article's own macro definitions that the retained text needs (\texttt{\textbackslash NN}, \texttt{\textbackslash cc}, \texttt{\textbackslash cont}, \texttt{\textbackslash ge}, \texttt{\textbackslash le}, \texttt{\textbackslash pmod}), with extra wrapper packages (bm, bbm, dsfont, xspace, cancel, mathtools). The delivered numbering is the unit's own. Assets: no retained span contains a figure environment, a graphics-inclusion command, a PDF-page inclusion, a table or a TikZ picture; no image file is copied into this package, the package declares no asset dependency and no image file is redistributed with it. Licence: version arXiv:2608.10702v1 of this article is distributed under the Creative Commons Attribution 4.0 International licence, as recorded in the arXiv licence metadata for this exact version and shown on the abstract page https://arxiv.org/abs/2608.10702v1. A full rights notice is delivered at licenses/arxiv-2608.10702-CC-BY-4.0.txt. Characters: every retained excerpt is pure ASCII, so the delivered engine, XeLaTeX with its default Latin Modern text font, reports no missing character.
\begin{lemma}\label{lem:squarefull-kernel}
Let $A:\NN\to\{\pm1\}$ be multiplicative, with $A(p)=1$ for every prime $p\ge5$ and $A(p^r)=1$ for $p\in\{2,3\}$ and $r\ge1$.
For $z\in[5,\infty]$ and $n\in\NN$ define
\[
A(n;z)=\sum_{\substack{d\mid n\\P^+(d)\le z}}(A*\mu)(d),\quad
C_A(u,q;z)=\sum_{\substack{d\ge1,\ P^+(d)\le z\\\gcd(d,q)\mid u}}
(A*\mu)(d)\frac{\gcd(d,q)}d.
\]
Extend $A(\,\cdot\,;z)$ evenly to nonzero integers and put $A(0;z)=0$.
Then
\begin{enumerate}
\item $A*\mu$ is multiplicative, supported on squarefull integers coprime to $6$, and $|(A*\mu)(p^r)|\le2$;
\item for every $\varepsilon>0$,
\[
\sum_{\substack{1\le n\le N\\n\equiv u\pmod q}}A(n;z)
=C_A(u,q;z)\frac Nq+O(N^{1/2+\varepsilon}),
\]
uniformly in $N,q,u,z$, and $|C_A(u,q;z)|\le1$;
\item if $q\mid Q$, every prime dividing $Q/q$ belongs to $\{2,3\}$, and $v\equiv u\pmod q$, then $C_A(v,Q;z)=C_A(u,q;z)$.
\end{enumerate}
\end{lemma}

\begin{lemma}\label{lem:nu-bounds}
Assume that $g_\cc$ is separable of degree $d\ge 2$.
\begin{enumerate}
\item If $v_p(\cont(g_\cc))\ge k$ then $\nu_\cc(p^k)=p^k$.
\item If $v_p(\cont(g_\cc))<k$ then $\nu_\cc(p^k)\le d\,p^{k(1-1/d)+v_p(\Delta(g_\cc))/d}$.
\item If $p\nmid \cont(g_\cc)$, then $\nu_\cc(p^k)\le 2p^{v_p(\Delta(g_\cc))/2}+d-2$.
\item If $p\nmid \Delta(g_\cc)$, then $\nu_\cc(p^k)=\nu_\cc(p)\le d$.
\end{enumerate}
\end{lemma}

\begin{proposition}\label{prop:squarefull-polynomial-values}
Let $A:\NN\to\{\pm1\}$ be multiplicative, with $A(p)=1$ for every prime $p\ge5$ and $A(p^r)=1$ for $p\in\{2,3\}$ and $r\ge1$.
Put
\begin{equation}\label{eq:def-KAp}
h_A=A*\mu,\quad
K_{A,p}(\cc)=1+\sum_{r\ge2}h_A(p^r)\xi_\cc(p^r),
\end{equation}
and let $\cc\in\mathcal G_{\mathrm{arith}}(H)$.
Then, for every $\varepsilon>0$,
\[
\sum_{1\le n\le x}A\bigl(g_\cc(n);\eta\bigr)
=\Bigl(\prod_{5\le p\le\eta}K_{A,p}(\cc)\Bigr)x
+O(x^\varepsilon),
\]
uniformly in $A$ and $\cc$.
\end{proposition}

\begin{proof}
Put
\begin{equation}\label{eq:height-cutoff}
Z=2dHx^d\asymp Hx^d=x^{5d}.
\end{equation}
Omit the at most $d$ integers $n$ for which $g_\cc(n)=0$.
On the remaining terms the convolution formula in Lemma~\ref{lem:squarefull-kernel} gives
\[
A(g_\cc(n);\eta)
=\sum_{\substack{q\mid g_\cc(n)\\P^+(q)\le\eta}}h_A(q).
\]
Including the integer roots temporarily in the congruence counts and then subtracting their contribution gives
\begin{align}\label{eq:trunc-main-err}
\sum_{1\le n\le x}A(g_\cc(n);\eta)
&=x\sum_{\substack{q\ge1\\P^+(q)\le\eta}}
\frac{h_A(q)\nu_\cc(q)}q-\mathcal T_A(Z;\cc)
+O\bigl(\mathcal M_A(Z;\cc)+x^\varepsilon\bigr),
\end{align}
where
\begin{equation}\label{eq:def-trunc-errors}
\begin{aligned}
\mathcal T_A(Z;\cc)
&=x\sum_{\substack{q>Z\\P^+(q)\le\eta}}
\frac{h_A(q)\nu_\cc(q)}q,\quad
\mathcal M_A(Z;\cc)=
\sum_{\substack{q\le Z\\P^+(q)\le\eta}}
|h_A(q)|\nu_\cc(q).
\end{aligned}
\end{equation}
Indeed, before the integer roots are removed, periodicity gives $x\nu_\cc(q)/q+O(\nu_\cc(q))$.
The Euler product $\prod_{p\le\eta}(1+2\lfloor\log Z/\log p\rfloor)$ gives
\[
\sum_{\substack{q\le Z\\P^+(q)\le\eta}}|h_A(q)|
\ll x^\varepsilon,
\]
so the integer roots contribute $O(x^\varepsilon)$.

Since $h_A$ and $\nu_\cc$ are multiplicative, Lemma~\ref{lem:nu-bounds}(2) shows that the local series below are absolutely convergent.
Hence
\begin{equation}\label{eq:trunc-euler-prod}
\sum_{\substack{q\ge1\\P^+(q)\le\eta}}
\frac{h_A(q)\nu_\cc(q)}q
=\prod_{p\le\eta}
\left(1+\sum_{r\ge2}h_A(p^r)\xi_\cc(p^r)\right)
=\prod_{5\le p\le\eta}K_{A,p}(\cc).
\end{equation}
The last equality uses $h_A(p^r)=0$ for $p\in\{2,3\}$.

We next record the uniform bound for $\nu_\cc(p^r)$ used in both error terms.
Fix $p\le\eta$ and put $v=v_p(\cont(g_\cc))$.
Since $\cont(g_\cc)\le\theta$, we have $p^v\le\theta$.
If $r\le v$, then $\nu_\cc(p^r)=p^r\le\theta$.
If $r>v$, write $g_\cc=p^v g_1$ and use Lemma~\ref{lem:nu-bounds}(3), together with
\[
\nu_\cc(p^r)=p^v\nu_{g_1}(p^{r-v}),\quad
\Delta(g_\cc)=p^{v(2d-2)}\Delta(g_1),\quad
v_p(\Delta(g_\cc))\le\theta.
\]
It follows that, for every $r\ge0$,
\begin{equation}\label{eq:nu-smallp-bound}
\nu_\cc(p^r)\ll\theta p^{\theta/2}.
\end{equation}

Let $R_p=\lfloor\log Z/\log p\rfloor$.
The bounds $h_A(p)=0$ and $|h_A(p^r)|\le2$ give
\[
\mathcal M_A(Z;\cc)
\le\prod_{p\le\eta}
\left(1+O\bigl(\theta R_p p^{\theta/2}\bigr)\right).
\]
The logarithm of the product is
\[
\ll\sum_{p\le\eta}\log(\theta R_p)
+\frac{\theta}{2}\sum_{p\le\eta}\log p
\ll\eta+\theta\eta
\ll\frac{\log x}{\log\log\log x}.
\]
Thus, uniformly in $A$ and $\cc$,
\begin{equation}\label{eq:E2-bound}
\mathcal M_A(Z;\cc)\ll x^\varepsilon.
\end{equation}

Finally, Rankin's trick gives
\begin{equation}\label{eq:large-prime-tail-rankin}
|\mathcal T_A(Z;\cc)|
\le xZ^{-1/2}\prod_{p\le\eta}
\left(1+\sum_{r\ge2}
\frac{|h_A(p^r)|\nu_\cc(p^r)}{p^{r/2}}\right).
\end{equation}
By \eqref{eq:nu-smallp-bound}, each inner sum is $O(\theta p^{\theta/2-1})$; hence the product is $O(x^\varepsilon)$ by the same logarithmic estimate.
Since $Z\asymp x^{5d}$, we obtain
\begin{equation}\label{eq:large-prime-tail-bound}
|\mathcal T_A(Z;\cc)|\ll1.
\end{equation}
Substitution in \eqref{eq:trunc-main-err} proves the proposition.
\end{proof}

\end{document}
